% GENERATED FILE. Edit canonical lessons, then run npm run generate:books.
% canonical-source-hash: fnv1a64:d2689799961d939b
% canonical-lessons: JA-C130-masshiro, JA-C130-makkuro, JA-C130-makka, JA-C130-aimai, JA-C130-tashika, JA-R130-first-pass-describing-words, JA-R130-second-pass-describing-people-and-colours

\chapter{Pure white, Pitch black, Bright red, Vague, Certain}
\label{ch:ja-pure-white-pitch-black-bright-red-vague-certain}
\hlchaptermodality{130}

\begin{chapteropening}
\textbf{By the end of this chapter:} \emph{I can say pure white, pitch black, bright red, vague and certain.}

Complete the last lesson of chapter 130: I can say pure white, pitch black, bright red, vague and certain.
\end{chapteropening}

\section[\texorpdfstring{\ja{まっしろ} (masshiro)}{まっしろ (masshiro)}]{\ja{まっしろ} (masshiro) — pure white}

\label{lesson:JA-C130-masshiro}



\begin{quote}
Before the new one: say the Japanese for shameless, then the Japanese for reliable.
\end{quote}

\subsection*{You'll want to know: \ja{まっしろ}}
\textbf{\ja{まっしろ}} — \emph{masshiro} — \textquotedblleft{}pure white\textquotedblright{}.

\begin{grammarlens}[title={What you've built}]
One more word for this chapter.
\end{grammarlens}

\subsection*{Your turn}
\begin{itemize}
\raggedright
  \item \emph{Say it:} \emph{masshiro}
  \item \emph{Say it:} \emph{masshiro}, once more
  \item \emph{Say it:} say \emph{masshiro}
  \item \emph{Recall:} say the Japanese for shameless, then the Japanese for reliable, then say \emph{masshiro} again
\end{itemize}

\begin{tcolorbox}[breakable,colback=teal!4,colframe=teal!35!black,title={Before you move on}]
What does \ja{まっしろ} mean? (\textquotedblleft{}Pure white\textquotedblright{}.) Say it once more.
\end{tcolorbox}
\section[\texorpdfstring{\ja{まっくろ} (makkuro)}{まっくろ (makkuro)}]{\ja{まっくろ} (makkuro) — pitch black}

\label{lesson:JA-C130-makkuro}



\begin{quote}
Before the new one: say the Japanese for reliable, then the Japanese for pure white.
\end{quote}

\subsection*{You'll want to know: \ja{まっくろ}}
\textbf{\ja{まっくろ}} — \emph{makkuro} — \textquotedblleft{}pitch black\textquotedblright{}.

\begin{grammarlens}[title={What you've built}]
One more word for this chapter.
\end{grammarlens}

\subsection*{Your turn}
\begin{itemize}
\raggedright
  \item \emph{Say it:} \emph{makkuro}
  \item \emph{Say it:} \emph{makkuro}, once more
  \item \emph{Say it:} say \emph{makkuro}
  \item \emph{Recall:} say the Japanese for reliable, then the Japanese for pure white, then say \emph{makkuro} again
\end{itemize}

\begin{tcolorbox}[breakable,colback=teal!4,colframe=teal!35!black,title={Before you move on}]
What does \ja{まっくろ} mean? (\textquotedblleft{}Pitch black\textquotedblright{}.) Say it once more.
\end{tcolorbox}
\section[\texorpdfstring{\ja{まっか} (makka)}{まっか (makka)}]{\ja{まっか} (makka) — bright red}

\label{lesson:JA-C130-makka}



\begin{quote}
Before the new one: say the Japanese for pure white, then the Japanese for pitch black.
\end{quote}

\subsection*{You'll want to know: \ja{まっか}}
\textbf{\ja{まっか}} — \emph{makka} — \textquotedblleft{}bright red\textquotedblright{}.

\begin{grammarlens}[title={What you've built}]
One more word for this chapter.
\end{grammarlens}

\subsection*{Your turn}
\begin{itemize}
\raggedright
  \item \emph{Say it:} \emph{makka}
  \item \emph{Say it:} \emph{makka}, once more
  \item \emph{Say it:} say \emph{makka}
  \item \emph{Recall:} say the Japanese for pure white, then the Japanese for pitch black, then say \emph{makka} again
\end{itemize}

\begin{tcolorbox}[breakable,colback=teal!4,colframe=teal!35!black,title={Before you move on}]
What does \ja{まっか} mean? (\textquotedblleft{}Bright red\textquotedblright{}.) Say it once more.
\end{tcolorbox}
\section[\texorpdfstring{\ja{あいまい} (aimai)}{あいまい (aimai)}]{\ja{あいまい} (aimai) — vague}

\label{lesson:JA-C130-aimai}



\begin{quote}
Before the new one: say the Japanese for pitch black, then the Japanese for bright red.
\end{quote}

\subsection*{You'll want to know: \ja{あいまい}}
\textbf{\ja{あいまい}} — \emph{aimai} — \textquotedblleft{}vague\textquotedblright{}.

\begin{grammarlens}[title={What you've built}]
One more word for this chapter.
\end{grammarlens}

\subsection*{Your turn}
\begin{itemize}
\raggedright
  \item \emph{Say it:} \emph{aimai}
  \item \emph{Say it:} \emph{aimai}, once more
  \item \emph{Say it:} say \emph{aimai}
  \item \emph{Recall:} say the Japanese for pitch black, then the Japanese for bright red, then say \emph{aimai} again
\end{itemize}

\begin{tcolorbox}[breakable,colback=teal!4,colframe=teal!35!black,title={Before you move on}]
What does \ja{あいまい} mean? (\textquotedblleft{}Vague\textquotedblright{}.) Say it once more.
\end{tcolorbox}
\section[\texorpdfstring{\ja{たしか} (tashika)}{たしか (tashika)}]{\ja{たしか} (tashika) — certain, sure}

\label{lesson:JA-C130-tashika}



\begin{quote}
Before the new one: say the Japanese for bright red, then the Japanese for vague.
\end{quote}

\subsection*{You'll want to know: \ja{たしか}}
\textbf{\ja{たしか}} — \emph{tashika} — \textquotedblleft{}certain, sure\textquotedblright{}.

\begin{grammarlens}[title={What you've built}]
One more word for this chapter.
\end{grammarlens}

\subsection*{Your turn}
\begin{itemize}
\raggedright
  \item \emph{Say it:} \emph{tashika}
  \item \emph{Say it:} \emph{tashika}, once more
  \item \emph{Say it:} say \emph{tashika}
  \item \emph{Recall:} say the Japanese for bright red, then the Japanese for vague, then say \emph{tashika} again
\end{itemize}

\begin{tcolorbox}[breakable,colback=teal!4,colframe=teal!35!black,title={Before you move on}]
What does \ja{たしか} mean? (\textquotedblleft{}Certain, sure\textquotedblright{}.) Say it once more.
\end{tcolorbox}
\section[(dialogue)]{Describing words — a first pass}

\label{lesson:JA-R130-first-pass-describing-words}



\begin{quote}
Say the words for vague and certain.
\end{quote}

\subsection*{Your turn}
Say these aloud:

\begin{itemize}
\raggedright
  \item square, pink, grey, thin, shallow
  \item deep, precious, necessary, simple, complicated
  \item well-behaved, dear, frustrating, suspicious, detailed
  \item itchy, smelly, humid, funny; strange, interesting
  \item amazing, terrible, young, very young, plentiful
\end{itemize}

\begin{tcolorbox}[breakable,colback=teal!4,colframe=teal!35!black,title={Before you move on}]
Square? (\textbf{\ja{しかくい}}, \emph{shikakui}.) Pink? (\textbf{\ja{ももいろ}}, \emph{momoiro}.) Grey? (\textbf{\ja{はいいろ}}, \emph{haiiro}.) Thin? (\textbf{\ja{うすい}}, \emph{usui}.) Shallow? (\textbf{\ja{あさい}}, \emph{asai}.) Deep? (\textbf{\ja{ふかい}}, \emph{fukai}.) Precious? (\textbf{\ja{たいせつ}}, \emph{taisetsu}.) Necessary? (\textbf{\ja{ひつよう}}, \emph{hitsuyou}.) Simple? (\textbf{\ja{かんたん}}, \emph{kantan}.) Complicated? (\textbf{\ja{ふくざつ}}, \emph{fukuzatsu}.) Well-behaved? (\textbf{\ja{おとなしい}}, \emph{otonashii}.) Dear? (\textbf{\ja{なつかしい}}, \emph{natsukashii}.) Frustrating? (\textbf{\ja{くやしい}}, \emph{kuyashii}.) Suspicious? (\textbf{\ja{あやしい}}, \emph{ayashii}.) Detailed? (\textbf{\ja{くわしい}}, \emph{kuwashii}.) Itchy? (\textbf{\ja{かゆい}}, \emph{kayui}.) Smelly? (\textbf{\ja{くさい}}, \emph{kusai}.) Humid? (\textbf{\ja{むしあつい}}, \emph{mushiatsui}.) Funny; strange? (\textbf{\ja{おかしい}}, \emph{okashii}.) Interesting? (\textbf{\ja{おもしろい}}, \emph{omoshiroi}.) Amazing? (\textbf{\ja{すごい}}, \emph{sugoi}.) Terrible? (\textbf{\ja{ひどい}}, \emph{hidoi}.) Young? (\textbf{\ja{わかい}}, \emph{wakai}.) Very young? (\textbf{\ja{おさない}}, \emph{osanai}.) Plentiful? (\textbf{\ja{ゆたか}}, \emph{yutaka}.)
\end{tcolorbox}
\section[(dialogue)]{Describing people and colours — a second pass}

\label{lesson:JA-R130-second-pass-describing-people-and-colours}



\begin{quote}
Say the words for happy and drowsy.
\end{quote}

\subsection*{Your turn}
Say these aloud:

\begin{itemize}
\raggedright
  \item happy, drowsy, relieved, ordinary, tasty; skilful
  \item many, few, refreshing, calm, honest
  \item clever, foolish, brave, fine, ugly
  \item cute, persistent, noisy, shameless, reliable
  \item pure white, pitch black, bright red, vague, certain
\end{itemize}

\begin{tcolorbox}[breakable,colback=teal!4,colframe=teal!35!black,title={Before you move on}]
Happy? (\textbf{\ja{しあわせ}}, \emph{shiawase}.) Drowsy? (\textbf{\ja{ねむたい}}, \emph{nemutai}.) Relieved? (\textbf{\ja{あんしん}}, \emph{anshin}.) Ordinary? (\textbf{\ja{ふつう}}, \emph{futsuu}.) Tasty; skilful? (\textbf{\ja{うまい}}, \emph{umai}.) Many? (\textbf{\ja{おおい}}, \emph{ooi}.) Few? (\textbf{\ja{すくない}}, \emph{sukunai}.) Refreshing? (\textbf{\ja{さわやか}}, \emph{sawayaka}.) Calm? (\textbf{\ja{おだやか}}, \emph{odayaka}.) Honest? (\textbf{\ja{すなお}}, \emph{sunao}.) Clever? (\textbf{\ja{かしこい}}, \emph{kashikoi}.) Foolish? (\textbf{\ja{おろか}}, \emph{oroka}.) Brave? (\textbf{\ja{いさましい}}, \emph{isamashii}.) Fine? (\textbf{\ja{こまかい}}, \emph{komakai}.) Ugly? (\textbf{\ja{みにくい}}, \emph{minikui}.) Cute? (\textbf{\ja{かわいい}}, \emph{kawaii}.) Persistent? (\textbf{\ja{しつこい}}, \emph{shitsukoi}.) Noisy? (\textbf{\ja{やかましい}}, \emph{yakamashii}.) Shameless? (\textbf{\ja{あつかましい}}, \emph{atsukamashii}.) Reliable? (\textbf{\ja{たのもしい}}, \emph{tanomoshii}.) Pure white? (\textbf{\ja{まっしろ}}, \emph{masshiro}.) Pitch black? (\textbf{\ja{まっくろ}}, \emph{makkuro}.) Bright red? (\textbf{\ja{まっか}}, \emph{makka}.) Vague? (\textbf{\ja{あいまい}}, \emph{aimai}.) Certain? (\textbf{\ja{たしか}}, \emph{tashika}.)
\end{tcolorbox}
